% source: https://tex.stackexchange.com/a/766815 (question 766814, score 5, block 1)
% Source - https://tex.stackexchange.com/q/766814
% Posted by minthao_2011
% Retrieved 2026-09-30, License - CC BY-SA 4.0
%https://tex.stackexchange.com/questions/766814/how-to-use-conditional-if-then-statements-in-luacas-lualatex-to-filter-algebra
\documentclass{article}
\usepackage{luacas}
\begin{document}
\begin{CAS}
vars("x")
left = x+5
right = x+5
f = topoly(right^2)
g = topoly(left)
eq = f-g
r = ZTable( roots(eq) )
sol1 = substitute({[x]=r[1]}, right):simplify()
sol2 = substitute({[x]=r[2]}, right):simplify()
sol1v = tostring(sol1) %expression of sol1 as a string
sol2v = tostring(sol2) %expression of sol1 as a string
\end{CAS}

Solve the equation:
\[\sqrt{\print{left}} = \print{right}.\]

Squaring both sides of the equation, we get:
\[\print{left} = \print{f}.\]

Simplifying the equation above yields:
\[ \print{eq} = 0. \] 

This equation has two roots:
\[r_1 = \print{r[1]}, \quad r_2 = \print{r[2]}.\]

\directlua{%
local sol, nb = {}, 0
if string.find(sol1v, "i") then
    tex.print("\\par These solution are not real.")
else
    tex.print("We have:")
    sol1v = load("return "..sol1v)() %numerical evaluation of sol1
    sol2v = load("return "..sol2v)() %numerical evaluation of sol2
    if sol1v >= 0 then 
        tex.print("\\[-r_1 + 5 = \\print{sol1} \\geq 0,\\] hence $r_1$ is a valid solution.")
        table.insert(sol,1); nb = nb+1
    else
        tex.print("\\[-r_1 + 5 = \\print{sol1} < 0,\\] hence $r_1$ is not a valid solution.")
    end
    if sol2v >= 0 then 
        tex.print("\\[-r_2 + 5 = \\print{sol2} \\geq 0,\\] hence $r_2$ is a valid solution.")
        table.insert(sol,2); nb = nb+1
    else
        tex.print("\\[-r_2 + 5 = \\print{sol2} < 0,\\] hence $r_2$ is not a valid solution.")
    end
    if nb == 0 then tex.print("\\par Therefore, there is no solution.")
    elseif nb == 1 then tex.print("\\par Therefore, there is one solution:") 
        if sol[1] == 1 then tex.print("$\\print{r[1]}$.") else tex.print("$\\print{r[1]}$.") end
    else 
        tex.print("\\par Therefore, there are two solution: $\\print{r[1]}$ and $\\print{r[2]}$.")
    end
end
}%
\end{document}
